% SGA 3, Expose IV, all of section 4.6.
% French authority: Exp4-13oct24.pdf, component-local pages 48--53
% (current-edition pages 224--229; original-print pages 220--227).
% Stops immediately before section 4.7.
% The external English Markdown was used only as a comparison.

\setcounter{footnote}{42}

\SGARefTarget{sga3:iv:subsection:4:6}\subsection*{4.6. 层关于等价关系的商的描述}
\addcontentsline{toc}{subsection}{\SGARefLink{sga3:iv:subsection:4:6}{4.6}. 层关于等价关系的商}

回顾假设, 给定的拓扑 \(\mathcal T\) 比典范拓扑更粗.

\medskip
\noindent\SGARefTarget{sga3:iv:proposition:4:6:1}\textbf{命题 4.6.1.}
\textit{设
\(
 R \mathrel{\substack{\xrightarrow{\ p_1\ }\\[-0.6ex]
                       \xrightarrow[\ p_2\ ]{}}} X
\)
是层 \(X\) 中的一个 \(\widetilde{\mathscr C}\)-等价关系.
如下定义 \(F\in\operatorname{Ob}(\widehat{\mathscr C})\):
对每个 \(S\in\operatorname{Ob}(\mathscr C)\),\footnote{编者注:
\(R\times S\) 表示 \(X\times S\) 中由
\(R\times S_{\mathrm{diag}}\subset
X\times X\times S\times S\) 定义的等价关系,
而 \(R_Z\) 表示它在 \(Z\) 上诱导的等价关系
(参见 \SGARefLink{sga3:iv:definition:3:1:6}{3.1.6}).}}
\[
 F(S)=\left\{\begin{array}{l}
 \text{\(S\)-层 \(Z\), 为 \(X_S\) 的子层且在 \(R\times S\) 下稳定,}\cr
 \text{其关于 \(R_Z\) 的商为 \(S\), 即 \(R_Z\rightrightarrows Z\to S\) 正合}
\end{array}\right\}.
\]
\textit{对每个层 \(Y\), \(\operatorname{Hom}(Y,F)\) 自然地等同于集合}
\[
 \left\{\begin{array}{l}
 \text{\(Y\)-层, 为 \(X\times Y\) 的子层且在 \(R\times Y\) 下稳定,}\cr
 \text{且其商为 \(Y\)}
 \end{array}\right\}.
\]
\textit{特别地, 子层 \(R\subset X\times X\) 对应于一个元素
\(p\in\operatorname{Hom}(X,F)\), 且正合图}
\[
 R\mathrel{\substack{\xrightarrow{\ p_1\ }\\[-0.6ex]
                     \xrightarrow[\ p_2\ ]{}}}X\xrightarrow{p}F
\]
\textit{将 \(F\) 等同于层商 \(X/R\).}

事实上, 置 \(Q=X/R\).
对每个层 \(Y\) 及每个态射 \(f\in\operatorname{Hom}(Y,Q)\),
它对应于一个截面 \(s:Y\hookrightarrow Q\times Y\), 考虑图
\[
\tag{*}\label{sga3:iv:proposition:4:6:1:eq:unnamed:01}
\begin{array}{ccccc}
 R\times Y&\rightrightarrows&X\times Y&\longrightarrow&Q\times Y\cr
 &&\big\uparrow&&\big\uparrow\!{\scriptstyle s}\cr
 &&Z&\longrightarrow&Y,
\end{array}
\]
其中右方块是笛卡尔方块.
由 \SGARefLink{sga3:iv:proposition:4:4:11}{4.4.11}, 立即可见
\(Z\) 是 \(X\times Y\) 的一个子 \(Y\)-层,
在 \(R\times Y\) 下稳定, 且商为 \(Y\).
反之, 每个这样的 \(Z\) 都来自 \(Q\times Y\) 在 \(Y\) 上的唯一截面.
先取 \(Y\) 为可表层, 得到同构 \(Q\simeq F\);
再取任意 \(Y\), 便得到所断言的 \(\operatorname{Hom}(Y,F)\) 的描述.
最后, 典范态射 \(X\to Q\) 对应于子 \(X\)-层
\(R\subset X\times X\), 证明完成.

\medskip
\noindent\SGARefTarget{sga3:iv:corollary:4:6:2}\textbf{推论 4.6.2.}
\textit{设 \(G\) 是 \(F\) 的任意子函子,
使得 \(\operatorname{Hom}(X,G)\subset\operatorname{Hom}(X,F)\)
包含 \(R\). 则典范态射 \(p:X\to F\) 通过 \(G\) 分解.}

由于 \(p\) 是覆盖态射
(\SGARefLink{sga3:iv:proposition:4:4:9}{4.4.9} 和
\SGARefLink{sga3:iv:proposition:4:4:3}{4.4.3}),
\(G\) 是 \(F\) 的一个加细.
特别地, \(F\) 的每个满足\SGARefLink{sga3:iv:corollary:4:6:2}{上述条件}
的子层 \(G\) 都等于 \(F\)
(\SGARefLink{sga3:iv:corollary:4:3:12}{4.3.12}).

\medskip
\noindent\SGARefTarget{sga3:iv:para:4:6:3}\textbf{4.6.3.}
现在考虑 \(X\) 与 \(R\) 可表的情形.
先引入一些术语. 除 \SGARefLink{sga3:iv:para:3:4:1}{3.4.1}
中引入的条件 \SGARefLink{sga3:iv:para:3:4:1:item:a:01}{(a)}--\SGARefLink{sga3:iv:para:3:4:1:item:d:01}{(d)} 外,
还将使用关于 \(\mathscr C\) 的态射族 \((M)\) 的下列条件;
为完整起见, 重述条件
\SGARefLink{sga3:iv:para:4:6:3:item:a:01}{(a)}--\SGARefLink{sga3:iv:para:4:6:3:item:c:01}{(c)}:
\begin{enumerate}
 \item[(a)]\SGARefTarget{sga3:iv:para:4:6:3:item:a:01}
 \((M)\) 在基变换下稳定.
 \item[(b)]\SGARefTarget{sga3:iv:para:4:6:3:item:b:01}
 \((M)\) 的两个元素的复合属于 \((M)\).
 \item[(c)]\SGARefTarget{sga3:iv:para:4:6:3:item:c:01}
 每个同构都属于 \((M)\).
 \item[\((d_{\mathcal T})\)]\SGARefTarget{sga3:iv:para:4:6:3:item:d:mathcal:t:01}
 \((M)\) 的每个元素都是覆盖态射.
 \footnote{编者注: \(\mathcal T\) 表示 \(\mathscr C\)
 上给定的拓扑, 它比典范拓扑更粗.}
 \item[\((e_{\mathcal T})\)]\SGARefTarget{sga3:iv:para:4:6:3:item:e:mathcal:t:01}
 设 \(f:X\to Y\) 是 \(\mathscr C\) 的一个态射.
 若存在 \(Y\) 的一个加细 \(R\),
 使得对每个 \(Y'\to R\), 态射 \(X\times_Y Y'\to Y'\)
 都属于 \((M)\), 则 \(f\) 属于 \((M)\).
\end{enumerate}
条件 \SGARefLink{sga3:iv:para:4:6:3:item:a:01}{(a)} 和
\SGARefLink{sga3:iv:para:4:6:3:item:b:01}{(b)} 蕴含
\begin{enumerate}
 \item[\((a')\)]\SGARefTarget{sga3:iv:para:4:6:3:item:aprime:01}
 \((M)\) 的两个元素的纤维积属于 \((M)\).
\end{enumerate}
另一方面, 由 \SGARefLink{sga3:iv:corollary:4:3:9}{4.3.9},
条件 \SGARefLink{sga3:iv:para:4:6:3:item:a:01}{(a)} 和
\SGARefLink{sga3:iv:para:4:6:3:item:d:mathcal:t:01}{\((d_{\mathcal T})\)} 蕴含
\begin{enumerate}
 \item[\((d')\)]\SGARefTarget{sga3:iv:para:4:6:3:item:dprime:01}
 \((M)\) 的每个元素都是普遍有效满态射.
\end{enumerate}

\medskip
\noindent\SGARefTarget{sga3:iv:para:4:6:4}\textbf{4.6.4.}
当 \(\mathscr C\) 有纤维积时, 记作 \((M_{\mathcal T})\) 的覆盖态射族
满足上述条件. 事实上 (参见 \SGARefLink{sga3:iv:definition:4:2:3}{4.2.3}),
\SGARefLink{sga3:iv:para:4:6:3:item:a:01}{(a)} 由
\SGARefLink{sga3:iv:definition:4:2:3:item:c:1:01}{(C~1)} 得到,
\SGARefLink{sga3:iv:para:4:6:3:item:b:01}{(b)} 由
\SGARefLink{sga3:iv:definition:4:2:3:item:c:2:01}{(C~2)} 得到,
\SGARefLink{sga3:iv:para:4:6:3:item:c:01}{(c)} 由
\SGARefLink{sga3:iv:definition:4:2:3:item:c:4:01}{(C~4)} 得到,
\SGARefLink{sga3:iv:para:4:6:3:item:d:mathcal:t:01}{\((d_{\mathcal T})\)} 由
\SGARefLink{sga3:iv:definition:4:2:3}{定义} 得到,
而 \SGARefLink{sga3:iv:para:4:6:3:item:e:mathcal:t:01}{\((e_{\mathcal T})\)} 由
\SGARefLink{sga3:iv:definition:4:2:3:item:c:5:01}{(C~5)} 得到.
所以, 下文对满足这些条件的族所建立的结果, 特别适用于
\((M_{\mathcal T})\). 特别地, 可以取 \(\mathcal T\) 为典范拓扑,
取 \((M)\) 为普遍有效满态射族.

\medskip
\noindent\SGARefTarget{sga3:iv:lemma:4:6:5}\textbf{引理 4.6.5.}
\textit{设 \((M)\) 是满足上述条件
\SGARefLink{sga3:iv:para:4:6:3:item:a:01}{(a)}--\SGARefLink{sga3:iv:para:4:6:3:item:e:mathcal:t:01}{\((e_{\mathcal T})\)}
的一个态射族.
设 \(R\) 是 \(X\in\operatorname{Ob}(\mathscr C)\) 中一个
\((M)\) 型的 \(\mathscr C\)-等价关系.
以 \(\widetilde X\) 表示 \(X\) 定义的层,
以 \(\widetilde R\) 表示 \(R\) 定义的 \(\widetilde X\) 中的
\(\widetilde{\mathscr C}\)-等价关系,
并以 \(\widetilde X/\widetilde R\) 表示层商.
则 \(R\) 是 \((M)\)-有效的, 当且仅当
\(\widetilde X/\widetilde R\) 可表.
这些条件成立时, \(\widetilde X/\widetilde R\) 由商 \(X/R\) 表出.}

先假设 \(R\) 是 \((M)\)-有效的, 并置 \(Y=X/R\).
典范态射 \(p:X\to Y\) 属于 \((M)\),
所以由 \SGARefLink{sga3:iv:para:4:6:3:item:d:mathcal:t:01}{\((d_{\mathcal T})\)},
它是覆盖态射. 因此相应的态射
\[
 \widetilde p:\widetilde X\longrightarrow
 \widetilde X/\widetilde R
\]
是 \(\widetilde{\mathscr C}\) 中的普遍有效满态射
(\SGARefLink{sga3:iv:proposition:4:4:3}{4.4.3}).
所以它将 \(\widetilde X/\widetilde R\) 等同于 \(\widetilde X\)
关于 \(\widetilde p\) 定义的等价关系 \(R'\) 的商.
由于典范函子 \(\mathscr C\to\widetilde{\mathscr C}\) 与纤维积可交换,
且 \(R\) 是 \(p\) 定义的等价关系, \(R'\) 恰好就是 \(\widetilde R\).

反之, 假设 \(\widetilde X/\widetilde R\) 由
\(\mathscr C\) 的一个对象 \(Y\) 表出.
设 \(p:X\to Y\) 是由典范态射
\(\widetilde X\to\widetilde X/\widetilde R\) 得到的态射;
由 \SGARefLink{sga3:iv:proposition:4:4:3}{4.4.3}, 它是覆盖态射.
如前, \(R\) 显然是 \(p\) 定义的等价关系.
只需再证明 \(p\in(M)\). 笛卡尔图
\[
\begin{array}{ccccc}
 R&\xrightarrow{\ \sim\ }&X\times_YX&\longrightarrow&X\cr
 {\scriptstyle p_2}\searrow&&\big\downarrow\!{\scriptstyle\mathrm{pr}_2}
 &&\big\downarrow\!{\scriptstyle p}\cr
 &&X&\xrightarrow{p}&Y
\end{array}
\]
表明, 在沿覆盖态射 \(p\) 作基变换后, 态射 \(p\) 变为
属于 \((M)\) 的 \(p_2\).
由 \SGARefLink{sga3:iv:para:4:6:3:item:e:mathcal:t:01}{\((e_{\mathcal T})\)}
即得\SGARefLink{sga3:iv:lemma:4:6:5}{此断言}.

\medskip
\noindent\SGARefTarget{sga3:iv:corollary:4:6:5:1}\textbf{推论 4.6.5.1.}
\footnote{编者注: \SGARefLink{sga3:iv:corollary:4:6:5:1}{本推论} 为新增内容.}
\textit{设 \((M)\) 满足上述条件
\SGARefLink{sga3:iv:para:4:6:3:item:a:01}{(a)}--\SGARefLink{sga3:iv:para:4:6:3:item:e:mathcal:t:01}{\((e_{\mathcal T})\)},
\(f:G\to G'\) 是一个满足 \(f\in(M)\) 的 \(\mathscr C\)-群态射.
假设 \(\operatorname{Ker}(f)\) 可表
(当 \(\mathscr C\) 有终对象 \(e\) 时即为如此).
则由 \(H=\operatorname{Ker}(f)\) 定义的 \(G\) 中的等价关系
是 \((M)\)-有效的,
且 \(G'\) 表出关于拓扑 \(\mathcal T\) 的层商
\(\widetilde G/\widetilde H\).}

这由 \SGARefLink{sga3:iv:lemma:4:6:5}{4.6.5} 和
\SGARefLink{sga3:iv:corollary:3:4:7:1}{3.4.7.1} 得到.

现在可以陈述本小节的主要结果.

\medskip
\noindent\SGARefTarget{sga3:iv:theorem:4:6:6}\textbf{定理 4.6.6.}
\textit{设 \((M)\) 是满足 \SGARefLink{sga3:iv:para:4:6:3}{4.6.3} 中公理
\SGARefLink{sga3:iv:para:4:6:3:item:a:01}{(a)}--\SGARefLink{sga3:iv:para:4:6:3:item:e:mathcal:t:01}{\((e_{\mathcal T})\)}
的一个态射族.
设 \(R\) 是 \(\mathscr C\) 的一个对象 \(X\) 中的
\((M)\) 型 \(\mathscr C\)-等价关系
(参见 \SGARefLink{sga3:iv:definition:3:4:3}{3.4.3}).
考虑由下式定义的函子 \(F\in\operatorname{Ob}(\widehat{\mathscr C})\):}
\[
 F(S)=\left\{\begin{array}{l}
 \text{\(S\)-层 \(Z\), 为 \(X_S\) 的子层且在 \(R\times S\) 下稳定,}\cr
 \text{其关于 \(R_Z\) 的商为 \(S\)}
 \end{array}\right\}.
\]
\textit{设 \(F_0\) 是如下定义的 \(F\) 的子函子:
\(F_0(S)\) 由可表的 \(Z\in F(S)\) 组成, 即}
\[
 F_0(S)=\left\{\begin{array}{l}
 \text{\(\mathscr C/S\)-对象 \(Z\), 为 \(X_S\) 的子对象且在 \(R\times S\) 下稳定, 并满足}\cr
 \text{\(R_Z\) 是 \((M)\)-有效的且商为 \(S\), 即满足}\cr
 \text{\(Z\to S\) 属于 \((M)\), 且
 \(R_Z\simeq Z\times_S Z\)}
 \end{array}\right\}.
\]
\textit{则:}
\begin{enumerate}
 \item[(i)]\SGARefTarget{sga3:iv:theorem:4:6:6:item:i:01}
 \textit{由子对象 \(R\subset X\times X\) 定义的态射
 \(p\in\operatorname{Hom}(X,F)=F(X)\),
 将 \(F\) 等同于 \(X\) 关于 \(R\) 的层商.}
 \item[(ii)]\SGARefTarget{sga3:iv:theorem:4:6:6:item:ii:01}
 \textit{以下条件等价:}
 \begin{enumerate}
  \item[(a)]\SGARefTarget{sga3:iv:theorem:4:6:6:item:a:01} \textit{\(F\) 可表;}
  \item[(b)]\SGARefTarget{sga3:iv:theorem:4:6:6:item:b:01} \textit{\(F_0\) 可表;}
  \item[(c)]\SGARefTarget{sga3:iv:theorem:4:6:6:item:c:01} \textit{\(R\) 是 \((M)\)-有效的.}
 \end{enumerate}
 \textit{在这些条件下, \(F=F_0=X/R\).}
 \item[(iii)]\SGARefTarget{sga3:iv:theorem:4:6:6:item:iii:01}
 \textit{设 \((N)\) 是一个在基变换下稳定的态射族,
 使得对每个覆盖族 \(\{S_i\to S\}\) 及每个属于 \((N)\)
 的态射族 \(\{T_i\to S_i\}\),
 \(T_i\) 上关于 \(\{S_i\to S\}\) 的每个下降数据都有效.
 假设 \(X\) 可平方 (参见 \SGARefLink{sga3:iv:definition:1:6:0}{1.6.0}),
 且 \(R\to X\times X\) 属于 \((N)\). 则 \(F_0=F\).}
\end{enumerate}

\noindent\textit{证明.}
断言 \SGARefLink{sga3:iv:theorem:4:6:6:item:i:01}{(i)}
已在 \SGARefLink{sga3:iv:proposition:4:6:1}{4.6.1} 中证明.

对于 \SGARefLink{sga3:iv:theorem:4:6:6:item:ii:01}{(ii)},
已经证明 \SGARefLink{sga3:iv:theorem:4:6:6:item:a:01}{(a)} 与
\SGARefLink{sga3:iv:theorem:4:6:6:item:c:01}{(c)} 等价, 以及等式 \(F=X/R\).
还需证明, \SGARefLink{sga3:iv:theorem:4:6:6:item:b:01}{(b)} 或
\SGARefLink{sga3:iv:theorem:4:6:6:item:c:01}{(c)} 都蕴含 \(F_0=F\).
首先注意, 如\SGARefLink{sga3:iv:theorem:4:6:6}{定理} 中所述,
\(F_0\) 确实是 \(F\) 的一个子函子:
若 \(S\in\operatorname{Ob}(\mathscr C)\) 且 \(Z\in F_0(S)\),
则 \(Z\to S\) 可平方, 所以对每个 \(S'\to S\), 有
\(Z\times_SS'\in F_0(S')\).
由于 \(R\in F(X)\) 属于 \(F_0(X)\),
由推论 \SGARefLink{sga3:iv:corollary:4:6:2}{4.6.2} 可知,
\SGARefLink{sga3:iv:theorem:4:6:6:item:b:01}{(b)} 蕴含 \(F_0=F\).

现在假设 \SGARefLink{sga3:iv:theorem:4:6:6:item:c:01}{(c)} 成立,
设 \(Q\) 是 \(\mathscr C\) 中表出 \(X/R\) 的一个对象.
则 \(X\to Q\) 属于 \((M)\).
对每个 \(S\in\operatorname{Ob}(\mathscr C)\) 及每个 \(Z\in F(S)\),
\SGARefLink{sga3:iv:proposition:4:6:1}{4.6.1} 中的图
\(\SGARefLink{sga3:iv:proposition:4:6:1:eq:unnamed:01}{(*)}\) 给出
\[
 Z=S\times_{Q\times S}(X\times S).
\]
所以 \(Z\) 可表, 且 \(Z\to S\) 属于 \((M)\), 因而 \(Z\in F_0(S)\).

对于 \SGARefLink{sga3:iv:theorem:4:6:6:item:iii:01}{(iii)},
设 \(f\in\operatorname{Hom}(S,F)\) 对应于 \(Z\in F(S)\).
必须证明 \(f\) 通过 \(F_0\) 分解,
或等价地, \(Z\) 可表.
当 \(f\) 通过 \(X\) 分解时, 这立即由\SGARefLink{sga3:iv:lemma:4:6:7}{下列引理} 得到.

\medskip
\noindent\SGARefTarget{sga3:iv:lemma:4:6:7}\textbf{引理 4.6.7.}
\textit{设 \(x_0\in X(S)\).
\(x_0\) 在 \(F(S)\) 中的像,
对应于由以下两个笛卡尔方块定义的子层 \(Z\subset X_S\):}
\[
\begin{array}{ccccc}
 X_S&\xrightarrow{\mathrm{id}_{X_S}\times x_0}&
 X_S\times_SX_S&\longrightarrow&X\times X\cr
 \big\uparrow&&\big\uparrow&&\big\uparrow\cr
 Z&\longrightarrow&R_S&\longrightarrow&R.
\end{array}
\]

这立即由态射 \(X\to F\) 的描述得到.

回到\SGARefLink{sga3:iv:theorem:4:6:6}{定理} 的证明,
先假设 \(f\) 通过 \(X\) 分解.
则 \(Z\) 可表; 此外, 由于 \(R\to X\times X\) 属于 \((N)\),
\(Z\to X_S\) 也属于该族.

一般情形中, \(f\) 未必通过 \(X\) 分解.
然而, 由于 \(X\to F\) 是覆盖态射
(\SGARefLink{sga3:iv:proposition:4:4:3}{4.4.3}),
由 \SGARefLink{sga3:iv:remark:4:4:8:item:vii:01}{4.4.8(vii)},
存在一个覆盖族 \(\{S_i\to S\}\),
以及对每个 \(i\), 一个使下图交换的态射 \(S_i\to X\):
\SGARefTarget{sga3:iv:supp:lemma:4:6:7:diagram:01}\[
\begin{array}{ccc}
 X&\longrightarrow&F\cr
 \big\uparrow&&\big\uparrow\!{\scriptstyle f}\cr
 S_i&\longrightarrow&S
\end{array}
\]
根据以上论证, \SGARefLink{sga3:iv:supp:lemma:4:6:7:diagram:01}{此图}
定义的态射 \(f_i:S_i\to F\) 属于 \(\operatorname{Hom}(S_i,F_0)\),
并对应于子层 \(Z\times_SS_i\subset X_{S_i}\).
态射 \(Z\times_SS_i\to X_{S_i}\) 属于 \((N)\),
且族 \(\{X_{S_i}\to X_S\}\) 是覆盖族.
所以只需再证明\SGARefLink{sga3:iv:proposition:4:6:8}{下列命题}.

\medskip
\noindent\SGARefTarget{sga3:iv:proposition:4:6:8}\textbf{命题 4.6.8.}
\textit{设 \(\{S_i\to S\}\) 是一个覆盖族,
\(Z\) 是 \(S\) 上的一个层.
假设对每个 \(i\), \(S_i\)-函子 \(Z\times_SS_i\) 由一个对象 \(T_i\) 表出.
则族 \((T_i)\) 带有关于 \(\{S_i\to S\}\) 的一个典范下降数据.
层 \(Z\) 可表当且仅当此数据有效;
此时, 下降所得的对象表出 \(Z\).}

首先注意, 由 \SGARefLink{sga3:iv:proposition:4:4:3}{4.4.3},
\(\{S_i\to S\}\) 是 \(\widetilde{\mathscr C}\) 中的一个普遍有效满态射族,
从而是 \(\widetilde{\mathscr C}\) 中的一个下降族
(\SGARefLink{sga3:iv:proposition:2:3}{2.3}).
若 \(Z\) 由一个对象 \(T\) 表出,
则 \(T\times_SS_i\) 视为层时同构于 \(Z\times_SS_i\).
因此, \(T_i\) 上的下降数据有效,
其唯一确定的下降对象同构于 \(T\), 因而表出 \(Z\).

反之, 假设 \(T_i\) 上的典范下降数据有效,
设 \(T\) 是下降所得的对象.
由于 \(\{S_i\to S\}\) 是 \(\widetilde{\mathscr C}\) 中的下降族,
存在一个 \(S\)-态射 \(T\to Z\),
其到每个 \(S_i\) 的基变换是典范态射 \(T_i\to Z\times_SS_i\).
此态射局部为同构; 由于 \(T\) 与 \(Z\) 都是层,
由 \SGARefLink{sga3:iv:corollary:4:5:8}{4.5.8} 可知它是同构.

\medskip
\noindent\SGARefTarget{sga3:iv:corollary:4:6:9}\textbf{推论 4.6.9.}
\textit{设 \(R\) 是 \(X\) 中的一个 \((M)\)-有效等价关系.
对每个层 \(F\), 映射}
\[
 \operatorname{Hom}(X/R,F)\longrightarrow\operatorname{Hom}(X,F)
\]
\textit{将第一个集合等同于第二个集合中与 \(R\) 相容的态射所成的子集.}

\medskip
\noindent\SGARefTarget{sga3:iv:corollary:4:6:10}\textbf{推论 4.6.10.}
\textit{设 \(\mathcal T'\) 是一个比 \(\mathcal T\) 更粗的拓扑,
且 \((M)\) 中的态射关于此拓扑都是覆盖态射.
在 \SGARefLink{sga3:iv:theorem:4:6:6:item:iii:01}{4.6.6(iii)} 的假设下,
对介于 \(\mathcal T'\) 和典范拓扑之间的每个拓扑,
\(X/R\) 也都是 \(X\) 关于 \(R\) 的层商.}

\medskip
\noindent\SGARefTarget{sga3:iv:remark:4:6:11}\textbf{注 4.6.11.}
若在 \SGARefLink{sga3:iv:theorem:4:6:6:item:iii:01}{4.6.6(iii)} 的陈述中,
再假设每当该处的假设满足时, 下降所得的态射 \(T\to S\) 都属于 \((N)\),
则包含态射 \(Z\hookrightarrow X_S\) 也属于 \((N)\).
这立即由 \(Z\) 的下降构造得到.

\medskip
\noindent\SGARefTarget{sga3:iv:remark:4:6:12}\textbf{注 4.6.12.}
蕴含关系
\[
 \text{\SGARefLink{sga3:iv:theorem:4:6:6:item:c:01}{(c)}}\Longrightarrow\text{\SGARefLink{sga3:iv:theorem:4:6:6:item:b:01}{(b)}}\Longrightarrow\text{\SGARefLink{sga3:iv:theorem:4:6:6:item:a:01}{(a)}}
 \qquad\text{及}\qquad
 \text{\SGARefLink{sga3:iv:theorem:4:6:6:item:c:01}{(c)}}\Longrightarrow[F_0=F=X/R]
\]
的证明没有使用引理 \SGARefLink{sga3:iv:lemma:4:6:5}{4.6.5} 的充分性方向,
而条件 \SGARefLink{sga3:iv:para:4:6:3:item:e:mathcal:t:01}{\((e_{\mathcal T})\)}
只在此方向中用到.
所以, 当 \((M)\) 仅满足条件
\SGARefLink{sga3:iv:para:4:6:3:item:a:01}{(a)}--\SGARefLink{sga3:iv:para:4:6:3:item:d:mathcal:t:01}{\((d_{\mathcal T})\)}
时, 这些蕴含仍成立.
可平方的覆盖态射族就是这样的族 \((M)\) 的一个例子
(比较 \SGARefLink{sga3:iv:para:4:6:4}{4.6.4}).
对典范拓扑, 这些态射恰好就是普遍有效满态射. 所以:

\medskip
\noindent\SGARefTarget{sga3:iv:corollary:4:6:13}\textbf{推论 4.6.13.}
\textit{设 \(R\) 是 \(X\) 中的一个普遍有效等价关系.
则 \(\mathscr C\) 的对象 \(X/R\) 是关于典范拓扑的
\(X\) 关于 \(R\) 的层商.
它表出下列函子: \((X/R)(S)\) 是 \(X_S\) 的那些在
\(R\times S\) 下稳定的子 \(\mathscr C/S\)-对象 \(Z\) 的集合,
其中诱导的等价关系普遍有效, 且商为 \(S\).}

类似地, 对任意拓扑:

\medskip
\noindent\SGARefTarget{sga3:iv:corollary:4:6:14}\textbf{推论 4.6.14.}
\textit{设 \((M)\) 是可平方的覆盖态射族.
若 \(R\) 是 \(X\) 中的一个 \((M)\)-有效等价关系,
则 \(\mathscr C\) 的对象 \(X/R\) 是 \(X\) 关于 \(R\) 的层商,
且表出 \SGARefLink{sga3:iv:theorem:4:6:6}{4.6.6} 中的函子 \(F_0\).}

\medskip
\noindent\SGARefTarget{sga3:iv:scholium:4:6:15}\textbf{附论 4.6.15.}
现在可以对 \SGARefLink{sga3:iv:scholium:4:5:5}{4.5.5} 作如下说明.
在仅涉及逆极限的问题中, 如纤维积、代数结构等,
第 I 篇的结果及 \SGARefLink{sga3:iv:scholium:4:5:5}{4.5.5} 允许我们
不加区别地将 \(\mathscr C\) 等同于 \(\widetilde{\mathscr C}\)
或 \(\widehat{\mathscr C}\) 的一个全子范畴.
对同时涉及逆极限和正向极限的问题, 则不能如此处理.

当逆极限和正向极限同时出现时, 特别是在取商时
(例如, 在一个群模其正规子群的商上赋予群结构时),
都将给定范畴视为嵌入层范畴中.
因此, 若 \(R\) 是 \(\mathscr C\) 的一个对象 \(X\) 中的
\(\mathscr C\)-等价关系, 则记号 \(X/R\) 表示 \(X\) 关于
\(R\) 的层商, 此前记作 \(j(X)/j(R)\);
而当此层可表时, 也表示表出它的对象.
以上结果表明, 在最重要的情形中, \(\mathscr C\) 中的商也是层范畴中的商.
无论如何, 若 \(\mathscr C\) 中的一个商与
\(\widetilde{\mathscr C}\) 中的商不一致,
例如不是普遍的商, 则不对它使用记号 \(X/R\).
这修改了\SGARefLink{sga3:iv:section:3}{第 3 号} 中的定义.

所以, 研究这类问题时, 先在层范畴中进行,
其中通常的结果都成立 (参见\SGARefLink{sga3:iv:subsection:4:4}{第 4.4 号});
随后利用本小节的结果, 以及在可用时利用下降有效性的判别法,
将所得陈述具体化到原来的范畴中.
以下小节将给出运用此方法的例子.

% Translation coverage: entire assigned source component, without omissions.
% Statistics: source lines=415; target lines=449; inline math=345; display math=13; unresolved translation TODOs=0.
